\chapter{Introduction}
%
\label{chap:Introduction}
%
\section{Motivation}
%
Modern industrial companies accumulate thousands of service tickets every year.
A single ticket at an agricultural machinery manufacturer such as Amazone~Werke
can span dozens of email exchanges, multi-page PDF invoices, spreadsheet attachments,
and scanned delivery notes—all in a mix of German, English, and occasionally French
or Dutch. Extracting actionable information from this heterogeneous corpus manually
is slow, inconsistent, and impossible to scale.

Advances in Large Language Models~(LLMs) offer a promising alternative. Models
such as Llama~3, Gemma~3, and GPT-4o have demonstrated strong zero-shot
capabilities in information extraction and text classification, even in specialised,
domain-specific settings. However, applying these models to real-world industrial
service data—with its multilingual content, verbose email threads, and
domain-specific abbreviations—requires a principled experimental evaluation that
is largely absent from the academic literature.

This thesis addresses that gap. It designs and evaluates an end-to-end pipeline that
(1)~converts raw ticket files into a unified textual representation, (2)~uses LLMs
to extract structured insights~(problem statement, solution, entity identifiers, and
categorical tags), and~(3)~rigorously benchmarks ten model configurations against
human-annotated ground truth.

%----------------------------------------------------------------------
\section{Problem Statement}
\label{sec:problem}
%
Service tickets at Amazone~Werke exhibit several properties that make automated
analysis challenging:

\begin{itemize}
    \item \textbf{Unstructured, heterogeneous content.} Important details are spread
          across email bodies~(MSG), PDF documents~(credit notes, inspection reports),
          Excel spreadsheets, and scanned images.  No consistent schema exists.

    \item \textbf{Multi-lingual text.} The same ticket may contain paragraphs in
          German, English, and technical abbreviations from both languages.
          Standard NLP tools optimised for a single language fail to process
          such code-switched content reliably.

    \item \textbf{Mixed structured and narrative information.} Document
          categories~(Gutschrift, Belastungsanzeige), article codes, and machine
          serial numbers follow predictable patterns that rule-based extraction
          handles well.  The underlying problem, root cause, and resolution require
          contextual reasoning that rules cannot provide.

    \item \textbf{Volume and scalability.} The corpus comprises tens of thousands
          of tickets.  Manual analysis is not feasible; any automated system must
          process tickets with minimal human intervention.

    \item \textbf{Incomplete written records.} Many ticket resolutions are agreed
          upon in phone calls or meetings, leaving the written record incomplete.
          Models must reason under uncertainty without hallucinating solutions.
\end{itemize}

%----------------------------------------------------------------------
\section{Research Questions}
\label{sec:rq}
%
The thesis is structured around one main research question and two sub-questions.

\begin{description}
    \item[\textbf{RQ~(Main):}] \emph{Can Large Language Models reliably extract
          actionable insights from industrial service tickets?}

    \item[\textbf{RQ\,1:}] \emph{Which ticket information can be captured reliably
          with rule-based extraction, and which parts require LLMs?}

    \item[\textbf{RQ\,2:}] \emph{Which LLM produces the most useful and coherent
          insights on the benchmark tickets?}
\end{description}

%----------------------------------------------------------------------
\section{Scope and Delimitations}
\label{sec:scope}
%
This thesis focuses on the \emph{extraction} of structured information from existing
ticket files; it does not cover automated ticket routing, priority prediction, or
resolution recommendation.  The evaluation is conducted on a benchmark of
250~tickets drawn from the Amazone~Werke service corpus; generalisation to other
industrial domains is discussed qualitatively in Chapter~\ref{chap:Conclusion}.

Models are evaluated in a zero-shot configuration using a fixed, domain-aware prompt.
Fine-tuning and retrieval-augmented generation~(RAG) are left as directions for
future work.

%----------------------------------------------------------------------
\section{Thesis Structure}
\label{sec:structure}
%
The remainder of this thesis is organised as follows.
Chapter~\ref{chap:RelatedWork} reviews related work on LLM-based information
extraction, service ticket analysis, and multi-label text classification.
Chapter~\ref{chap:Methodology} describes the overall methodology, including the
preprocessing pipeline, benchmark construction, and the experimental setup.
Chapter~\ref{chap:Results} presents the quantitative results and qualitative
observations.
Chapter~\ref{chap:Discussion} interprets the findings and discusses the
limitations of the approach.
Chapter~\ref{chap:Conclusion} summarises the findings, answers the research
questions, and outlines directions for future work.
